\begin{table}[ht]
\centering
\caption{Fiscal Multipliers, United Kingdom (AR(1) Fiscal Rules)}
\label{tab:mult_uk}
\begin{tabular}{lcccc}
\hline\hline
Instrument & $H=1$ & $H=4$ & $H=12$ & $H=20$ \\
\hline
Government spending      & 1.02 & 0.94 & 0.68 & 0.56 \\
                         & [0.77, 1.18] & [0.70, 1.09] & [0.52, 0.85] & [0.36, 0.78] \\
Labor tax                & 0.08 & 0.25 & 0.61 & 0.74 \\
                         & [0.03, 0.15] & [0.14, 0.39] & [0.46, 0.75] & [0.56, 0.94] \\
Capital tax              & 0.92 & 1.56 & 2.48 & 2.61 \\
                         & [0.73, 1.29] & [1.13, 2.45] & [1.78, 3.44] & [1.87, 3.52] \\
Consumption tax          & 0.30 & 0.38 & 0.38 & 0.32 \\
                         & [0.23, 0.40] & [0.32, 0.46] & [0.30, 0.43] & [0.22, 0.40] \\
Expensing rate           & 0.26 & 0.49 & 0.74 & 0.95 \\
                         & [0.03, 0.42] & [0.22, 0.70] & [0.47, 1.03] & [0.59, 1.34] \\
\hline\hline
\multicolumn{5}{p{0.85\textwidth}}{\footnotesize Note: Present-value own-revenue multipliers (equation~\ref{eq:mult_main}). Posterior median with 5--95 percentile intervals in brackets, from 1{,}000 posterior draws. A positive value means output moves in the expansionary direction of the policy. UK model with AR(1) fiscal rules, estimated on 1997Q2--2016Q4 (Appendix~\ref{app:uk}). The UK model has no investment tax credit.}
\end{tabular}
\end{table}
